\documentclass[../../main.tex]{subfiles}
\begin{document}
\begin{oframed}
  \subfile{problem}
\end{oframed}

{\bf [解答]}

与えられた数列
\begin{align}
  \begin{cases}
    a_{n+2}=a_{n+1}+a_n, a_0=1,a_1=2 \\
    b_n \equiv a_n \pmod 3 \quad (b_n=0,1,2) \\
    c_n=b_0+\cdots+b_n
  \end{cases} \label{eq:1}
\end{align}
について考える．以下$n \in \mathbb{Z}_{\ge 0}$ とする．

\paragraph{(1)}

\cref{eq:1}の漸化式から法を3として $a_{n+2} \equiv a_{n+1} + a_n$ だから，
\begin{align}
  &b_0 = 1,\ b_1 = 2,\ b_2 = 0,\ b_3 = 2,\ b_4 = 2,\ b_5 = 1,\ b_6 = 0, \\
  &b_7 = 1,\ b_8 = 1,\ b_9 = 2
\end{align}
である．

\paragraph{(2)}

$c_n$ の定義から
\begin{align}
  c_{n+8} - c_n = \sum_{k=1}^{8}b_{n+k} \label{eq:2}
\end{align}
である．
(1)から $\{b_n\}$ は $\{1, 2, 0, 2, 2, 1, 0, 1\}$ の周期数列だから，
\cref{eq:2}の右辺には$n$にかかわらずこれらの要素が1つずつ入っているので，
\begin{align}
  \sum_{k=1}^{8}b_{n+k} = \sum_{k=0}^7 b_k = c_7
\end{align}
だから\cref{eq:2}に代入して
\begin{align}
  c_{n+8} = c_n + c_7 \quad \square \label{eq:3}
\end{align}
である．

\paragraph{(3)}

(1)より$c_7 = 9$ だから\cref{eq:3}は
\begin{align}
  c_{n+8} = c_n + 9 \label{eq:4}
\end{align}
である．以下，題意の不等式
\begin{align}
  n+1 \le c_n \le \dfrac{3}{2}(n+1) \label{eq:5}
\end{align}
を数学的帰納法で示す．

まず$n=0\sim 7$ の時，具体的に$c_0$から$c_7$を計算すると
\begin{align}
  1 \le c_0 = 1 \le \dfrac{3}{2}  && 5 \le c_4 = 7 \le 9 \\
  2 \le c_1 = 3 \le \dfrac{9}{2}  && 6 \le c_5 = 8 \le \dfrac{21}{2} \\
  3 \le c_2 = 3 \le 6             && 7 \le c_6 = 8 \le 12 \\
  4 \le c_3 = 5 \le \dfrac{15}{2} && 8 \le c_7 = 9 \le 15
\end{align}
より\cref{eq:5}は成立する．

そこで$n=8k, 8k+1, \cdots, 8k+7$ ($k \in \mathbb{Z}_{\ge 0}$) での成立を仮定すると，\cref{eq:5}の両辺に$9$を足して\cref{eq:4}より
\begin{align}
  & n+1+9 \le c_{n+8} \le \dfrac{3}{2}(n+1)+9 \\
  \therefore
  & n+9 \le c_{n+8} \le \dfrac{3}{2}(n+7) \le \dfrac{3}{2}(n+9)
\end{align}
より $n$ で成立するなら $n+8$ でも成立するから，$n=8(k+1)\sim 8(k+1)+7$ でも成立する．

以上から数学的帰納法により\cref{eq:5}は示された． $\square$

\end{document}
